% !TeX program = xelatex
%==============================================================
%  Python金融数据分析 · 第7章 数据输入输出
%  基于《Python金融大数据分析（第2版）》第9章内容改编
%  授课对象：本科金融系学生（已学完第2-6章）
%  编译：xelatex chapter7.tex（编译两遍）
%==============================================================
\documentclass[9pt]{ctexbeamer}

%%%%-----导入宏包-----%%%%
\usepackage{style/shubeamer}   % SHU 主题（配色、帧标题、页脚、Python代码样式）
\usepackage{amsmath,amsfonts,amssymb,bm}
\usepackage{graphicx,hyperref,url}
\usepackage{subcaption}
\usepackage{booktabs}
\usepackage[super,square]{natbib}
\usepackage{wrapfig}          % 文字绕排
\usepackage{calligra}         % 结尾页花体字

%%%%-----设置字体-----%%%%
\setmainfont{CMU Serif}
\setsansfont{CMU Sans Serif}

%%%%-----常用自定义颜色（正文可直接使用）-----%%%%
\definecolor{nvidia}{RGB}{102,156,28}
\definecolor{red}{RGB}{184,13,73}
\definecolor{orange}{RGB}{242,151,36}
\definecolor{darkteal}{RGB}{43,106,108}
\definecolor{darkblue}{RGB}{4,37,58}

%%%%-----每节开始时自动显示当前节目录-----%%%%
\AtBeginSection[]
{
	\begin{frame}<beamer>
		\frametitle{\textbf{目录}}
		\textbf{\tableofcontents[currentsection]}
	\end{frame}
}

%%%%-----设定图片存放路径-----%%%%
\graphicspath{{figures/}}

%%%%-----首页信息设置-----%%%%
\AtBeginDocument{%
	%%%%----短标题[显示在页脚]{封面标题}
	\title[Python金融数据分析 · 第7章]{\fontsize{16pt}{22pt}\selectfont {Python金融数据分析}}
	%%%%----副标题（本章主题）
	\subtitle{\fontsize{9pt}{14pt}\selectfont \textit{第7章\quad 数据输入输出}}
	%%%%----个人信息
	\author[王伟冠]{%
		\begin{tabular}{ll}%
			主讲人： & 王伟冠\tabularnewline%
			院\quad 系： & 经济学院金融系\tabularnewline
		\end{tabular}
	}
	%%%%----机构信息
	\institute[SHU]{上海大学}
	%%%%----日期信息
	\date[\today]{\today}}

%%%%-----数学宏（向量、矩阵、算子等，见 style/macros.tex）-----%%%%
\input{style/macros.tex}


\begin{document}

%% 封面
\begin{frame}
	\titlepage
\end{frame}

%% 目录页
\section*{目录}
\begin{frame}
	\frametitle{\textbf{目录}}
	\textbf{\tableofcontents}
\end{frame}


%==============================================================
\section{数据从哪里来，到哪里去}
%==============================================================

\begin{frame}{本章学习目标}
	\begin{itemize}
		\item 理解"分析结果要能保存、下次能复用"的工作流意识
		\item 会用 Python 内置功能读写文本文件（\texttt{open} + \texttt{with}）
		\item 熟练掌握 CSV：金融数据最常见的交换格式
		\item 了解 Excel 读写的写法与依赖库安装
		\item 初步认识 HDF5（PyTables）：大表格的高速二进制存储
		\item 会根据场景选择合适的数据格式
	\end{itemize}
	\begin{block}{本章数据}
		继续用第6章的真实数据 \texttt{eod\_data.csv}，形成"读入$\to$计算$\to$保存"的完整闭环。
	\end{block}
\end{frame}

\begin{frame}{金融分析中的数据流}
	\begin{itemize}
		\item 真实工作的典型流程：
		\begin{enumerate}
			\item 从数据商/交易所/数据库\textbf{读入}原始数据；
			\item 清洗、计算（前几章的技能）；
			\item 把结果\textbf{保存}下来：给报告用、给同事用、给明天的自己用。
		\end{enumerate}
		\item 常见格式一览：
	\end{itemize}
	\begin{table}
		\centering
		\begin{tabular}{lll}
			\toprule
			格式 & 特点 & 典型用途 \\
			\midrule
			CSV & 纯文本、通用、可读 & 行情、财报交换 \\
			Excel & 多工作表、人人会开 & 报告、手工整理 \\
			HDF5 & 二进制、快、省空间 & 大规模数值数据 \\
			\bottomrule
		\end{tabular}
	\end{table}
\end{frame}


%==============================================================
\section{Python基本文件操作}
%==============================================================

\begin{frame}[fragile]{open + with：读写文本}
	\begin{lstlisting}
# 写文件：把分析结论记下来
with open("notes.txt", "w", encoding="utf-8") as f:
    f.write("第6章结论：SPY存在明显厚尾\n")
    f.write("beta(AAPL) = 0.96\n")

# 读文件：逐行看回来
with open("notes.txt", encoding="utf-8") as f:
    for line in f:
        print(line.strip())
\end{lstlisting}
	\begin{itemize}
		\item \texttt{"w"} 写（覆盖）、\texttt{"r"} 读（默认）、\texttt{"a"} 追加。
		\item \texttt{with} 语句：用完自动关闭文件，不怕忘写 \texttt{close()}。
		\item 中文内容务必加 \texttt{encoding="utf-8"}，否则 Windows 上易乱码。
	\end{itemize}
	\begin{alertblock}{零基础提醒}
		\texttt{"w"} 会\alert{清空}原有内容！想追加请用 \texttt{"a"}。
	\end{alertblock}
\end{frame}


%==============================================================
\section{CSV：金融数据的通用语言}
%==============================================================

\begin{frame}[fragile]{读CSV：三件套参数}
	\begin{lstlisting}
import pandas as pd

df = pd.read_csv("eod_data.csv",
                 index_col=0,       # 第0列作索引
                 parse_dates=True)  # 索引解析为日期
df.head(3)
\end{lstlisting}
	\begin{itemize}
		\item CSV = Comma-Separated Values，逗号分隔的纯文本，用记事本都能打开。
		\item 第6章已见过：真实行情、财报数据绝大多数以 CSV 交付。
		\item 常用参数还有：\texttt{sep}（分隔符）、\texttt{skiprows}（跳过行）、
			\texttt{na\_values}（自定义缺失标记）。
	\end{itemize}
\end{frame}

\begin{frame}[fragile]{写CSV：把计算结果保存下来}
	\begin{lstlisting}
import numpy as np
rets = np.log(df / df.shift(1))

stats = pd.DataFrame({
    "年化收益": rets.mean() * 252,
    "年化波动率": rets.std() * np.sqrt(252)})
stats.to_csv("stats.csv")            # 默认带索引列
stats.to_csv("stats_no_index.csv",
             index=False)            # 不写索引
\end{lstlisting}
	\begin{itemize}
		\item \texttt{to\_csv} 与 \texttt{read\_csv} 是一对"镜像"操作。
		\item 下次打开 notebook，直接 \texttt{read\_csv("stats.csv")}——不用重算。
	\end{itemize}
	\begin{alertblock}{往返一致性}
		写出再读回，\texttt{shape}、数值应完全一致。养成"写后立即读回抽查"的习惯。
	\end{alertblock}
\end{frame}

\begin{frame}[fragile]{读回验证}
	\begin{lstlisting}
back = pd.read_csv("stats.csv", index_col=0)
back.head(3)
back.equals(stats)   # True：一来一回数据没变
\end{lstlisting}
	\begin{itemize}
		\item \texttt{equals()} 检查两个 DataFrame 是否完全相同。
		\item 注意：CSV 是文本格式，浮点数有打印精度限制；
			对精度极度敏感的场景用二进制格式（后面讲）。
	\end{itemize}
\end{frame}


%==============================================================
\section{Excel：报告和协作的标配}
%==============================================================

\begin{frame}[fragile]{读写Excel}
	\begin{lstlisting}
# 需要先安装引擎库（命令行执行一次即可）：
#   pip install openpyxl

df.to_excel("data.xlsx",            # 写
            sheet_name="价格")
stats.to_excel("data.xlsx",
               sheet_name="统计")   # 第二个表

pd.read_excel("data.xlsx",          # 读
              sheet_name="统计",
              index_col=0)
\end{lstlisting}
	\begin{itemize}
		\item pandas 本身不带 Excel 引擎，\texttt{.xlsx} 依赖 \texttt{openpyxl}。
		\item 一个文件可有多个 \textbf{sheet}（工作表），用 \texttt{sheet\_name} 指定。
		\item Excel 适合给人看；给程序用、数据量大时，首选 CSV 或 HDF5。
	\end{itemize}
\end{frame}

\begin{frame}{CSV vs Excel：怎么选？}
	\begin{itemize}
		\item \textbf{CSV}：
		\begin{itemize}
			\item 优点：小、快、通用（任何软件都能读）、适合版本管理；
			\item 缺点：只有一张表、无格式、无公式。
		\end{itemize}
		\item \textbf{Excel}：
		\begin{itemize}
			\item 优点：多工作表、可加颜色/公式、同事都熟悉；
			\item 缺点：文件大、读写慢、大表（百万行）会卡。
		\end{itemize}
	\end{itemize}
	\begin{block}{经验法则}
		中间数据、程序间交换 $\to$ CSV/HDF5；
		最终交付给业务同事的报告 $\to$ Excel。
	\end{block}
\end{frame}


%==============================================================
\section{HDF5：大数据的高速通道}
%==============================================================

\begin{frame}[fragile]{HDF5 读写：两行搞定}
	\begin{lstlisting}
# 写：complevel 开启压缩（引擎是PyTables，pip安装过tables即可）
df.to_hdf("data.h5", key="prices",
          mode="w", complevel=9, complib="zlib")

# 读：一行读回完整DataFrame，日期索引原样保留
back = pd.read_hdf("data.h5", "prices")
back.equals(df)     # True：一来一回数据没变
\end{lstlisting}
	\begin{itemize}
		\item HDF5 是\textbf{二进制}格式：读写快、可压缩，适合大规模数值数据。
		\item \texttt{key} 相当于文件里的"表名"，一个 	exttt{.h5} 文件可存多张表。
		\item 还支持\alert{部分读取}：百万行里只取满足条件的子集，不必全部载入内存。
	\end{itemize}
\end{frame}

\begin{frame}{格式对比：同一份数据的三种存法}
	\begin{table}
		\centering
		\begin{tabular}{lcl}
			\toprule
			格式 & 实测大小 & 特点 \\
			\midrule
			CSV & 约 203 KB & 通用但解析慢 \\
			HDF5（压缩） & 约 152 KB & 更小更快，程序友好 \\
			Excel & 通常最大 & 人能直接看，大表会卡 \\
			\bottomrule
		\end{tabular}
	\end{table}
	\begin{itemize}
		\item 同样是第6章的行情表（2138行$\times$12列）实测所得。
		\item 数据量越大，二进制格式的优势越明显——真实研究动辄千万行。
	\end{itemize}
	\begin{alertblock}{注意}
		HDF5 文件不能用记事本/Excel 打开，交换给他人前要确认对方能读。
	\end{alertblock}
\end{frame}


%==============================================================
\section{综合案例与小结}
%==============================================================

\begin{frame}[fragile]{综合案例：一条完整的分析流水线}
	\begin{lstlisting}
# 1. 读入原始数据（第6章已熟悉）
df = pd.read_csv("eod_data.csv", index_col=0,
                 parse_dates=True).dropna(subset=["SPY"])

# 2. 计算：对数收益率与滚动波动率
rets = np.log(df / df.shift(1))
vol = rets["SPY"].rolling(21).std() * np.sqrt(252)

# 3. 保存中间结果：明天不用重算
vol.dropna().to_csv("spy_vol21.csv", header=True)

# 4. 保存给报告的Excel版本
vol.dropna().to_excel("spy_vol21.xlsx")
\end{lstlisting}
	\begin{itemize}
		\item 读入 $\to$ 计算 $\to$ 保存：本章把前六章的技能串成真实工作流。
	\end{itemize}
\end{frame}

\begin{frame}{本章小结}
	\begin{itemize}
		\item \textbf{文本}：\texttt{open + with}，注意 \texttt{"w"} 覆盖与 \texttt{encoding}。
		\item \textbf{CSV}：\texttt{read\_csv / to\_csv} 镜像操作；
			三件套 \texttt{index\_col}、\texttt{parse\_dates}、\texttt{index=False}。
		\item \textbf{Excel}：\texttt{read\_excel / to\_excel} + \texttt{openpyxl}；
			多工作表用 \texttt{sheet\_name}。
		\item \textbf{HDF5}：\texttt{to\_hdf / read\_hdf}，二进制高速存储，大数据首选。
		\item \textbf{选型}：程序交换用 CSV/HDF5，人工报告用 Excel；
			写完立即读回验证是好习惯。
	\end{itemize}
\end{frame}

\begin{frame}{课后任务与下章预告}
	\begin{block}{课后任务}
		\begin{enumerate}
			\item 把第6章练习中算出的 MSFT beta 与 AAPL beta 写入一个 CSV 文件，
				再读回并打印。
			\item 用 \texttt{with open} 写一个文本文件，记录你对"厚尾现象"的三句话总结。
			\item 选做：对 \texttt{stats.csv} 分别用 \texttt{index=True} 与
				\texttt{index=False} 写出，用记事本打开比较两个文件的差别。
		\end{enumerate}
	\end{block}
	\begin{exampleblock}{下章预告}
		数据读进来了，计算却太慢？下一章讲 \textbf{性能}：
		为什么循环不如向量化，以及如何衡量代码的运行时间。
	\end{exampleblock}
	\begin{center}
		\vspace{0.5em}
		本章内容改编自教材第9章\citep{hilpisch2018python}。
	\end{center}
\end{frame}


%==============================================================
%  附录：参考文献 + 结尾页
%==============================================================
\appendix

\begin{frame}[allowframebreaks]{参考文献}
	\bibliographystyle{style/gbt7714-2005}
	\bibliography{reference.bib}
\end{frame}

\begin{frame}
	\begin{center}
		{\Huge\calligra Thanks for your attention!}
		\vspace{1cm}

		{\Huge Q \& A}
	\end{center}
\end{frame}

\end{document}
